\section{Two quadratic masses: the exact second derivative and collision calculus}
\label{hqj:sec:collision}

Let $P(y)=\sum_{d=0}^Dp_dy^d\in\C[y]$ and $u,v>0$. Define
\begin{align}
\mathcal Z_P(s;a;u,v)
&=\sum_{n\in\Z}P((n+a)^2)
\bigl(((n+a)^2+u)((n+a)^2+v)\bigr)^{-s},
\label{hqj:eq:product-spectrum}\\
E_{P,u}(s;a)
&=\sum_{n\in\Z}P((n+a)^2)((n+a)^2+u)^{-s}.
\label{hqj:eq:weighted-E}
\end{align}
Their initial absolute-convergence conditions are
$\Rea s>(2D+1)/4$ and $\Rea s>D+1/2$, respectively. We use their
meromorphic continuations below. The finite polynomial identity
\begin{equation}\label{hqj:eq:polynomial-translate}
E_{P,u}(s;a)=\sum_{\ell=0}^D
\frac{P^{(\ell)}(-u)}{\ell!}E(s-\ell;a,\sqrt u)
\end{equation}
will reduce the individual weighted jets to the already evaluated family.

\subsection{The complete subtraction, not a formal divergent logarithm sum}
For $r\ge2$, put
\begin{equation}\label{hqj:eq:Tr}
T_r(u,v)=(-1)^r\sum_{k=1}^{r-1}
\frac{(u^k-v^k)(u^{r-k}-v^{r-k})}{k(r-k)}.
\end{equation}
These are the coefficients of $y^{-r}$ in
$\log^2((y+u)/(y+v))$ at infinity. The first three are
\begin{align}
T_2&=(u-v)^2,\notag\\
T_3&=-(u-v)^2(u+v),\label{hqj:eq:Tlow}\\
T_4&=(u-v)^2\left\{\tfrac23(u^2+uv+v^2)+\tfrac14(u+v)^2\right\}.
\notag
\end{align}
The full polynomial part of the weighted logarithmic square is
\begin{equation}\label{hqj:eq:PiP}
\Pi_P(y;u,v)=\sum_{d=2}^Dp_d\sum_{r=2}^dT_r(u,v)y^{d-r}.
\end{equation}
Define the ordinary sum
\begin{equation}\label{hqj:eq:KP}
\mathcal K_P(a;u,v)=\sum_{n\in\Z}\left\{
P(x^2)\log^2\frac{x^2+u}{x^2+v}-\Pi_P(x^2;u,v)\right\},
\qquad x=n+a.
\end{equation}
Its summand is $O(x^{-2})$, locally uniformly in positive $u,v$.
For degree at most one no subtraction is needed. At larger degree the
subtraction is generally nonzero and must not be discarded.

\begin{theorem}[Second derivative of a two-mass spectrum]
\label{hqj:thm:product}
Both spectra in \eqref{hqj:eq:product-spectrum}--\eqref{hqj:eq:weighted-E}
are holomorphic at $s=0$. For every polynomial $P$,
\begin{align}
\mathcal Z_P''(0;a;u,v)
={}&2\sum_{\ell=0}^D\frac1{\ell!}
\bigl\{P^{(\ell)}(-u)M_\ell(a,\sqrt u)
+P^{(\ell)}(-v)M_\ell(a,\sqrt v)\bigr\}
\notag\\
&-\mathcal K_P(a;u,v).
\label{hqj:eq:product-second}
\end{align}
When $u=v$, the correction vanishes and the formula reduces to
$\mathcal Z_P''(0;a;u,u)=4E_{P,u}''(0;a)$, as required by
$\mathcal Z_P(s;a;u,u)=E_{P,u}(2s;a)$.
\end{theorem}
\begin{proof}
On either positive lattice tail, all mass expansions and the weight
$P(x^2)$ contain only even powers of $x$. A pole of the associated
Hurwitz factor at $s=0$ would require a power $x^{-1}$, which is
impossible. Finite-head subtractions therefore prove holomorphy at
zero exactly as in Section~\ref{hqj:sec:local}, without any residue
term.

Compare $\mathcal Z_P(s)$ with
$\tfrac12(E_{P,u}(2s)+E_{P,v}(2s))$. After taking out $x^{-4s}$,
the pointwise difference has zero constant and linear coefficients,
and its quadratic coefficient is
\[
\tfrac12P(x^2)\log^2\frac{x^2+u}{x^2+v}.
\]
There is no $x^{-1}$ residue to lower the spectral order. Thus its
second derivative is the continued value of the logarithmic-square
sum at the fixed lattice coordinate.

The only nonsummable terms of that sum are the polynomial
\eqref{hqj:eq:PiP}. Their bilateral regularized sums vanish exactly:
for every $k\ge0$,
\begin{equation}\label{hqj:eq:even-polynomial-zero}
\zeta(-2k,a)+\zeta(-2k,1-a)=0.
\end{equation}
For example, this follows from
$\zeta(-2k,a)=-B_{2k+1}(a)/(2k+1)$ and
$B_{2k+1}(1-a)=-B_{2k+1}(a)$, including $k=0$.
Consequently the regularized value is exactly the ordinary remainder
\eqref{hqj:eq:KP}, not that remainder plus an unknown constant.
Taking the second derivative of the comparison and using
\eqref{hqj:eq:polynomial-translate} proves
\eqref{hqj:eq:product-second}.
\end{proof}

For example, at $P(y)=y^2$ the full correction is
\[
\mathcal K_{y^2}(a;u,v)=\sum_{n\in\Z}
\left\{x^4\log^2\frac{x^2+u}{x^2+v}-(u-v)^2\right\}.
\]
At $P(y)=y^3$ it is
\[
\sum_{n\in\Z}\left\{x^6\log^2\frac{x^2+u}{x^2+v}
-(u-v)^2x^2+(u-v)^2(u+v)\right\}.
\]
These illustrate why positivity of an unsubtracted logarithmic square
does not imply positivity of every polynomially regularized correction.

\subsection{An ordinary one-dimensional polylogarithmic reduction}
The positive first resolvent has the elementary formula
\begin{equation}\label{hqj:eq:E1}
E(1;a,\sqrt t)
=\frac\pi{\sqrt t}
\frac{\sinh(2\pi\sqrt t)}{\cosh(2\pi\sqrt t)-\cos(2\pi a)}.
\end{equation}
One proof takes the difference of the cotangent partial fractions at
$a\pm i\sqrt t$; another differentiates the known $L_0$ in the squared
mass. Put
\begin{equation}\label{hqj:eq:RP}
\mathcal R_P(t)=P(-t)E(1;a,\sqrt t).
\end{equation}

\begin{theorem}[The two-mass integral and mixed derivative]
\label{hqj:thm:Kintegral}
For $u>v>0$,
\begin{equation}\label{hqj:eq:Kintegral}
\boxed{\mathcal K_P(a;u,v)=2\int_v^u
\mathcal R_P(t)\log\frac{u-t}{t-v}\dd t.}
\end{equation}
The endpoint singularities are integrable logarithms. The mixed
parameter derivative is
\begin{equation}\label{hqj:eq:Kmixed}
\partial_u\partial_v\mathcal K_P(a;u,v)
=-2\frac{\mathcal R_P(u)-\mathcal R_P(v)}{v-u},
\end{equation}
with the continuous coincident value $2\mathcal R_P'(u)$.
Writing $u=b^2$, $v=c^2$, $b>c>0$, the integral is equivalently
\begin{align}
\mathcal K_P(a;b^2,c^2)=4\pi\int_c^b P(-y^2)
&\left\{1+2\Rea\Li_0(e^{-2\pi y+2\pi ia})\right\}
\notag\\[-2pt]
&\times\log\frac{b^2-y^2}{y^2-c^2}\dd y.
\label{hqj:eq:Kpolylog}
\end{align}
For $u<v$ use symmetry; for $u=v$ the correction is zero.
\end{theorem}
\begin{proof}
Polynomial division and \eqref{hqj:eq:even-polynomial-zero} show that
the regularized weighted resolvent is
\[
\Reg\sum_{n\in\Z}\frac{P(x^2)}{x^2+t}
=P(-t)E(1;a,\sqrt t)=\mathcal R_P(t).
\]
The convention here is precisely polynomial subtraction in $x^2$.
Partial fractions therefore give
\begin{equation}\label{hqj:eq:double-resolvent}
\Reg\sum_{n\in\Z}\frac{P(x^2)}{(x^2+t)(x^2+w)}
=\frac{\mathcal R_P(t)-\mathcal R_P(w)}{w-t}.
\end{equation}
Its diagonal is removable. Differentiating
\eqref{hqj:eq:KP} with respect to $u,v$ is valid after subtracting the
corresponding polynomial coefficients; the remainder and its
parameter derivatives remain $O(x^{-2})$ on compact sets. Since the
mixed derivative of the logarithmic square is
$-2/((x^2+u)(x^2+v))$, this proves \eqref{hqj:eq:Kmixed}.

Alternatively write the logarithm as
$\int_v^u(x^2+t)^{-1}\dd t$ and square it. Polynomial subtraction
commutes with these integrals coefficient by coefficient; the remaining
uniformly convergent sum may be integrated by ordinary Fubini. Thus
$\mathcal K_P$ is the double integral of
\eqref{hqj:eq:double-resolvent} over $[v,u]^2$.
The same symmetric divided-difference calculation used in
Theorem~\ref{hqj:thm:gamma-square}, with its sign reversed, gives
\eqref{hqj:eq:Kintegral}.

Finally substitute $t=y^2$ and use the absolutely convergent Poisson
kernel identity
\[
\frac{\sinh(2\pi y)}{\cosh(2\pi y)-\cos(2\pi a)}
=1+2\Rea\frac{e^{-2\pi y+2\pi ia}}{1-e^{-2\pi y+2\pi ia}}
=1+2\Rea\Li_0(e^{-2\pi y+2\pi ia}).
\]
This proves \eqref{hqj:eq:Kpolylog} and its convergence.
\end{proof}
The formulas give the full second product derivative in terms of
finite negative-integer second jets and a displayed ordinary integral.
There is no implicit subtraction of two divergent logarithmic sums.

\subsection{Every collision coefficient is a finite polygamma expression}
\begin{theorem}[Convergent collision expansion]\label{hqj:thm:collision}
Let $g=\operatorname{dist}(a,\Z)=\min(a,1-a)$. For $u>0$ and
$|h|<u+g^2$, the correction has the convergent Taylor expansion
\begin{align}
\mathcal K_P(a;u+h,u)
={}&2\sum_{r\ge2}\frac{(-1)^rH_{r-1}}r h^r
\sum_{\ell=0}^{\min(D,r-1)}
\frac{P^{(\ell)}(-u)}{\ell!}
E(r-\ell;a,\sqrt u).
\label{hqj:eq:collision-series}
\end{align}
It defines the analytic branch continued from $h=0$ in that disk.
In particular,
\begin{equation}\label{hqj:eq:confluent}
\lim_{h\to0}\frac{\mathcal K_P(a;u+h,u)}{h^2}
=P(-u)E(2;a,\sqrt u)+P'(-u)E(1;a,\sqrt u)
=-\mathcal R_P'(u).
\end{equation}
All resolvents on the right of \eqref{hqj:eq:collision-series} have
positive integer order and admit finite polygamma expressions.
\end{theorem}
\begin{proof}
The elementary identity
\[
\log^2(1+z)=2\sum_{r\ge2}\frac{(-1)^rH_{r-1}}r z^r,
\qquad |z|<1,
\]
follows either by multiplying the logarithm series or by differentiating
both sides. Put $z=h/(x^2+u)$. Its modulus is uniformly less than one
on every compact subdisk $|h|<u+g^2$. The polynomial subtraction in
\eqref{hqj:eq:KP} depends holomorphically on $h$, and its remainder is
normally summable on such compact sets. Taylor coefficients may
therefore be extracted after summation and subtraction.

At order $r$, polynomial expansion around $-u$ gives
\[
\frac{P(x^2)}{(x^2+u)^r}
=\sum_{\ell=0}^D\frac{P^{(\ell)}(-u)}{\ell!}(x^2+u)^{\ell-r}.
\]
For $\ell\ge r$ the terms are polynomials in $x^2$ and have zero
bilateral regularized sum. For $\ell<r$ they are ordinarily summable
positive-order resolvents. This gives exactly
\eqref{hqj:eq:collision-series}. The $r=2$ coefficient yields the
first equality in \eqref{hqj:eq:confluent}; differentiating
\eqref{hqj:eq:RP} and using $\partial_uE(1)=-E(2)$ gives the second.
\end{proof}

For clarity, a finite polygamma formula can be written without a
recursive mass derivative. For $N\ge1$, define
\begin{align}
A_{N,r}(b)&=(-1)^{N-r}\binom{2N-r-1}{N-r}(-2ib)^{-(2N-r)},
\quad 1\le r\le N,\notag\\
\Xi_1(z)&=-\psi(z),\qquad
\Xi_r(z)=\frac{(-1)^r}{(r-1)!}\psi^{(r-1)}(z)\quad(r\ge2).
\end{align}
Then the elementary partial fractions at the conjugate roots give
\begin{equation}\label{hqj:eq:finite-polygamma}
E(N;a,b)=2\Rea\sum_{r=1}^N A_{N,r}(b)
\bigl\{\Xi_r(a+ib)+\Xi_r(1-a+ib)\bigr\}.
\end{equation}
For $r=1$ the two conjugate harmonic tails cancel before the digamma
values are taken. The formula can also be checked from
\begin{equation}\label{hqj:eq:resolvent-derivative}
E(N;a,\sqrt u)=\frac{(-1)^{N-1}}{(N-1)!}
\partial_u^{N-1}E(1;a,\sqrt u).
\end{equation}
Thus no noninteger-order or unevaluated infinite spectral sum is
required in any collision coefficient.

For $P=1$, the expansion begins
\begin{align}
\mathcal K_1(a;u+h,u)
={}&h^2E(2)-h^3E(3)+\tfrac{11}{12}h^4E(4)
-\tfrac56h^5E(5)+\cdots,
\label{hqj:eq:collision-low}
\end{align}
all at $(a,\sqrt u)$. For higher-degree $P$, the extra translated
polynomial derivatives in \eqref{hqj:eq:collision-series} are essential;
keeping only $P(-u)E(r)$ would be incorrect.

\subsection{An exact averaged logarithmic-square identity}
An explicit elementary consequence provides both an independent check
and a useful integral evaluation.
\begin{corollary}[Average over the lattice shift]\label{hqj:cor:mean}
For $b,c>0$,
\begin{align}
\int_0^1\mathcal K_1(a;b^2,c^2)\dd a
&=\int_{-\infty}^{\infty}
\log^2\frac{x^2+b^2}{x^2+c^2}\dd x\notag\\
&=8\pi\bigl\{b\log(2b)+c\log(2c)
-(b+c)\log(b+c)\bigr\}.
\label{hqj:eq:mean-K}
\end{align}
Consequently
\begin{align}
\int_0^1\mathcal Z_1''(0;a;b^2,c^2)\dd a
=8\pi\bigl\{&b\log(2b)+c\log(2c)
+(b+c)\log(b+c)\notag\\
&-2(b+c)\bigr\}.
\label{hqj:eq:mean-product}
\end{align}
\end{corollary}
\begin{proof}
Absolute convergence unfolds the shift average of the sum into the
real-line integral. Alternatively the mean of $E(1;a,\sqrt t)$ is
$\pi/\sqrt t$. Averaging \eqref{hqj:eq:Kintegral} and substituting
$t=y^2$ gives, for $b>c$,
\[
4\pi\int_c^b\log\frac{b^2-y^2}{y^2-c^2}\dd y.
\]
Split the logarithm into its four linear factors. The integrals of
$\log(b-y)$ and $-\log(y-c)$ cancel. Integrating the remaining
$\log(b+y)-\log(y+c)$ gives twice the braces in
\eqref{hqj:eq:mean-K}. Symmetry covers $c>b$ and continuity covers
$b=c$. The mean of \eqref{hqj:eq:M0} is its zero Fourier mode
$8\pi b(\log(2b)-1)$, so \eqref{hqj:eq:product-second} proves
\eqref{hqj:eq:mean-product}.
\end{proof}
